% generated by scripts/export-catalog.ts — engine v2, seed 42
% requires \usepackage[T1]{fontenc}, \usepackage[utf8]{inputenc}, \usepackage[italian]{babel}

\section{Disposizioni}

\subsection{\texorpdfstring{\texttt{disposizioni-semplici-codice} (base)}{disposizioni-semplici-codice (base)}}

\paragraph{Quesito.} Quanti codici di 3 cifre tutte diverse si possono formare usando le cifre da 1 a 8?

\paragraph{Risposta corretta.} $336$

\paragraph{Derivazione.} L'ordine conta e le cifre non si ripetono: sono disposizioni semplici $D(8,3) = 8 \cdot  7 \cdot  6 = 336$. Se le cifre potessero ripetersi sarebbero $8^{3} = 512$; se l'ordine non contasse sarebbero $C(8,3) = 56$.

\paragraph{Distrattori.}
\begin{enumerate}
  \item $512$ --- \texttt{with-repetition} (\texttt{repetition}): Hai consentito di ripetere le cifre: \guillemotleft{}tutte diverse\guillemotright{} esclude i codici con cifre ripetute, quindi si contano disposizioni semplici, non disposizioni con ripetizione ($8^{3}$).
  \item $56$ --- \texttt{order-dropped} (\texttt{order-irrelevant}): Hai ignorato l'ordine: un codice è una sequenza ordinata, quindi scambiare due cifre produce un codice diverso.
  \item $40320$ --- \texttt{all-digits-permuted} (\texttt{constraint-ignored}): Hai contato le permutazioni di tutte le 8 cifre: il codice usa solo 3 posizioni, quindi il prodotto si arresta dopo 3 fattori.
\end{enumerate}

\subsection{\texorpdfstring{\texttt{disposizioni-ripetizione-targa} (base)}{disposizioni-ripetizione-targa (base)}}

\paragraph{Quesito.} Un lucchetto ha 4 caselle e su ogni casella si possono impostare 6 simboli diversi, anche ripetuti. Quante combinazioni diverse esistono?

\paragraph{Risposta corretta.} $1296$

\paragraph{Derivazione.} Ogni casella è indipendente e i simboli si possono ripetere: $6^{4} = 1296$. 360 sarebbe il risultato se i simboli dovessero essere tutti diversi, mentre 4096 scambia la base con l'esponente.

\paragraph{Distrattori.}
\begin{enumerate}
  \item $360$ --- \texttt{repetition-excluded} (\texttt{repetition}): Hai escluso le ripetizioni: il testo dice esplicitamente che i simboli si possono ripetere su caselle diverse.
  \item $4096$ --- \texttt{base-exponent-swapped} (\texttt{other}): Hai scambiato la base con l'esponente: sono 4 caselle con 6 simboli ciascuna, quindi la base è 6 e l'esponente è 4.
  \item $15$ --- \texttt{order-dropped} (\texttt{order-irrelevant}): Hai trattato il codice come un insieme senza ordine: invece ogni posizione del lucchetto è distinta, quindi l'ordine conta.
\end{enumerate}

\subsection{\texorpdfstring{\texttt{disposizioni-podio} (intermedio)}{disposizioni-podio (intermedio)}}

\paragraph{Quesito.} A una gara partecipano 6 atleti. In quanti modi diversi possono essere assegnate le medaglie d'oro, d'argento e di bronzo?

\paragraph{Risposta corretta.} $120$

\paragraph{Derivazione.} Le tre medaglie sono diverse, quindi l'ordine conta: $D(6,3) = 6 \cdot  5 \cdot  4 = 120$. $C(6,3) = 20$ conterebbe solo i gruppi di tre atleti, senza distinguere chi vince l'oro.

\paragraph{Distrattori.}
\begin{enumerate}
  \item $20$ --- \texttt{order-dropped} (\texttt{order-irrelevant}): Hai contato solo i gruppi di tre atleti: il podio distingue oro, argento e bronzo, quindi l'ordine conta e servono le disposizioni.
  \item $216$ --- \texttt{with-repetition} (\texttt{repetition}): Hai permesso che lo stesso atleta vincesse più medaglie: chi prende l'oro non può più prendere l'argento, quindi i fattori calano a ogni assegnazione.
  \item $18$ --- \texttt{athletes-times-medals} (\texttt{other}): Hai moltiplicato gli atleti per le medaglie ($6 \cdot  3$): a ogni assegnazione gli atleti disponibili diminuiscono e i fattori vanno moltiplicati fra loro ($6 \cdot  5 \cdot  4$).
\end{enumerate}

\subsection{\texorpdfstring{\texttt{disposizioni-vincolo-posizione} (intermedio)}{disposizioni-vincolo-posizione (intermedio)}}

\paragraph{Quesito.} Quanti codici di 4 cifre tutte diverse iniziano con una cifra dispari e terminano con una cifra pari?

\paragraph{Risposta corretta.} $1400$

\paragraph{Derivazione.} Scegli prima le estremità: ultima cifra pari (5 scelte), prima cifra dispari (5 scelte, tutti i dispari sono ancora disponibili). Restano 8 cifre per la seconda e 7 per la terza: $5 \cdot  5 \cdot  8 \cdot  7 = 1400$.

\paragraph{Distrattori.}
\begin{enumerate}
  \item $5040$ --- \texttt{constraint-ignored} (\texttt{constraint-ignored}): Hai contato tutti i codici con cifre distinte, ignorando entrambi i vincoli di parità sulle estremità.
  \item $10000$ --- \texttt{with-repetition} (\texttt{repetition}): Hai permesso di ripetere le cifre: \guillemotleft{}tutte diverse\guillemotright{} richiede di escludere a ogni passo le cifre già usate.
  \item $1800$ --- \texttt{reading-order} (\texttt{other}): Hai scelto le posizioni nell'ordine di lettura: ma dopo le due cifre centrali le pari disponibili non sono sempre 5. Conviene fissare prima le due estremità ($5 \cdot  5$) e poi riempire il centro ($8 \cdot  7$).
\end{enumerate}

\section{Permutazioni}

\subsection{\texorpdfstring{\texttt{permutazioni-semplici-anagrammi} (base)}{permutazioni-semplici-anagrammi (base)}}

\paragraph{Quesito.} Quanti anagrammi, anche privi di significato, si possono formare con le lettere della parola CAMPO?

\paragraph{Risposta corretta.} $120$

\paragraph{Derivazione.} La parola CAMPO ha 5 lettere tutte diverse: gli anagrammi sono le permutazioni di 5 elementi, cioè $5! = 120$.

\paragraph{Distrattori.}
\begin{enumerate}
  \item $24$ --- \texttt{one-letter-fixed} (\texttt{other}): Hai fissato una lettera in una posizione fissa: nessuna lettera è vincolata, quindi tutte le 5 vanno permutate.
  \item $20$ --- \texttt{two-factors-only} (\texttt{other}): Hai arrestato il prodotto a due fattori: ogni lettera occupa una posizione, quindi il prodotto arriva a 5 fattori.
  \item $32$ --- \texttt{two-choices-per-letter} (\texttt{other}): Hai dato solo due alternative per lettera: in realtà ogni posizione può ospitare qualunque lettera non ancora usata.
\end{enumerate}

\subsection{\texorpdfstring{\texttt{permutazioni-ripetizione-anagrammi} (intermedio)}{permutazioni-ripetizione-anagrammi (intermedio)}}

\paragraph{Quesito.} Quanti anagrammi distinti, anche privi di significato, si possono formare con le lettere della parola SASSO?

\paragraph{Risposta corretta.} $20$

\paragraph{Derivazione.} SASSO ha 5 lettere con ripetizioni ($S\times 3$). Gli anagrammi distinti sono $5!$ / ($3!$) = 20. Usare solo $5! = 120$ conterebbe più volte lo stesso anagramma, perché scambiare fra loro lettere uguali non produce parole diverse.

\paragraph{Distrattori.}
\begin{enumerate}
  \item $120$ --- \texttt{identical-letters-distinguished} (\texttt{labelled-groups}): Hai contato più volte lo stesso anagramma: scambiare fra loro le lettere uguali (qui $S\times 3$) non produce parole diverse.
  \item $40$ --- \texttt{divided-by-multiplicity} (\texttt{other}): Hai diviso per la molteplicità (3) invece che per il suo fattoriale: le lettere uguali si scambiano fra loro in $3!$ modi.
  \item $10$ --- \texttt{positions-only} (\texttt{other}): Hai scelto soltanto le posizioni della lettera ripetuta, $C(5,3)$, senza disporre le lettere rimanenti nei posti liberi.
\end{enumerate}

\subsection{\texorpdfstring{\texttt{permutazioni-vincolo-vicinanza} (avanzato)}{permutazioni-vincolo-vicinanza (avanzato)}}

\paragraph{Quesito.} 7 amici si mettono in fila per una foto. In quanti modi possono disporsi se due di loro, Anna e Bruno, vogliono stare sempre uno accanto all'altro?

\paragraph{Risposta corretta.} $1440$

\paragraph{Derivazione.} Si considerano Anna e Bruno come un unico blocco: restano 6 elementi da permutare, cioè $6! = 720$ modi. Dentro il blocco i due possono scambiarsi fra loro in 2 modi, quindi $2 \cdot  6! = 1440$. Dimenticare lo scambio interno porta a 720; ignorare il vincolo porta a $7! = 5040$.

\paragraph{Distrattori.}
\begin{enumerate}
  \item $720$ --- \texttt{block-swap-forgotten} (\texttt{constraint-applied-once}): Hai dimenticato lo scambio interno al blocco: Anna e Bruno possono permutarsi fra loro in 2 modi.
  \item $5040$ --- \texttt{constraint-ignored} (\texttt{constraint-ignored}): Hai ignorato il vincolo: i due amici vanno trattati come un unico blocco, che riduce gli elementi da permutare a 6.
  \item $4320$ --- \texttt{subtracted-block} (\texttt{other}): Hai sottratto invece di ragionare sul blocco: il vincolo riduce gli elementi da permutare, non è un'esclusione di casi.
\end{enumerate}

\section{Combinazioni}

\subsection{\texorpdfstring{\texttt{combinazioni-semplici-commissione} (base)}{combinazioni-semplici-commissione (base)}}

\paragraph{Quesito.} In una classe di 9 studenti si deve scegliere una delegazione di 2 persone, senza assegnare incarichi diversi. Quante delegazioni diverse sono possibili?

\paragraph{Risposta corretta.} $36$

\paragraph{Derivazione.} Nessun incarico distingue i membri, quindi l'ordine non conta: $C(9,2) = 9!$ / ($2! \cdot  7!$) = 36. Se i ruoli fossero diversi l'ordine conterebbe e si otterrebbe $D(9,2) = 72$.

\paragraph{Distrattori.}
\begin{enumerate}
  \item $72$ --- \texttt{order-counted} (\texttt{order-irrelevant}): Hai trattato l'ordine come rilevante: la delegazione non ha incarichi, quindi ogni gruppo di persone è una sola scelta e servono le combinazioni.
  \item $18$ --- \texttt{students-times-seats} (\texttt{other}): Hai moltiplicato il numero di studenti per il numero di posti ($9 \cdot  2$): non esiste una regola che dia i gruppi come prodotto di questi due numeri; servono le combinazioni $C(9,2)$.
  \item $9$ --- \texttt{one-member-short} (\texttt{other}): Hai contato un membro di meno: la delegazione è di 2 persone, non 1.
\end{enumerate}

\subsection{\texorpdfstring{\texttt{combinazioni-sottoinsiemi} (avanzato)}{combinazioni-sottoinsiemi (avanzato)}}

\paragraph{Quesito.} Quanti sottoinsiemi, compresi l'insieme vuoto e l'insieme stesso, ha un insieme di 7 elementi?

\paragraph{Risposta corretta.} $128$

\paragraph{Derivazione.} Per ogni elemento si decide se includerlo o no: $2 \cdot  2$ $\cdot$ \ldots{} = $2^{7} = 128$. Equivalentemente è la somma di tutte le combinazioni $C(7,0)$ + \ldots{} + $C(7,7)$. Il valore 127 esclude l'insieme vuoto.

\paragraph{Distrattori.}
\begin{enumerate}
  \item $127$ --- \texttt{empty-set-excluded} (\texttt{constraint-ignored}): Hai escluso l'insieme vuoto: il testo chiede di comprenderlo insieme all'insieme stesso.
  \item $5040$ --- \texttt{orderings-counted} (\texttt{order-irrelevant}): Hai contato gli ordinamenti degli elementi: un sottoinsieme non ha ordine interno, quindi non servono le permutazioni.
  \item $21$ --- \texttt{pairs-only} (\texttt{other}): Hai contato solo i sottoinsiemi di due elementi: invece ogni cardinalità da 0 a 7 contribuisce al totale.
\end{enumerate}

\subsection{\texorpdfstring{\texttt{combinazioni-quote-gruppi} (intermedio)}{combinazioni-quote-gruppi (intermedio)}}

\paragraph{Quesito.} Un gruppo è formato da 7 ragazze e 5 ragazzi. In quanti modi si può comporre una squadra di 2 ragazze e 2 ragazzi?

\paragraph{Risposta corretta.} $210$

\paragraph{Derivazione.} Le due scelte sono indipendenti e vanno moltiplicate: $C(7,2) \cdot  C(5,2) = 21 \cdot  10 = 210$. Sommare i due valori (31) sarebbe corretto solo per scelte alternative, non simultanee; $C(12,4) = 495$ ignora le quote richieste.

\paragraph{Distrattori.}
\begin{enumerate}
  \item $31$ --- \texttt{added-choices} (\texttt{additive-union}): Hai sommato le due scelte: la somma vale per scelte alternative, ma qui le quote vanno soddisfatte contemporaneamente, quindi si moltiplicano.
  \item $495$ --- \texttt{quotas-ignored} (\texttt{constraint-ignored}): Hai unito i due gruppi ignorando le quote: la squadra deve avere esattamente 2 ragazze e 2 ragazzi, non 4 persone a caso.
  \item $840$ --- \texttt{order-counted} (\texttt{order-irrelevant}): Hai introdotto un ordine che non esiste: all'interno di ciascun gruppo le persone si scelgono senza incarichi, quindi servono le combinazioni.
\end{enumerate}

\subsection{\texorpdfstring{\texttt{combinazioni-diagonali-poligono} (intermedio)}{combinazioni-diagonali-poligono (intermedio)}}

\paragraph{Quesito.} Quante diagonali ha un poligono convesso con 10 vertici?

\paragraph{Risposta corretta.} $35$

\paragraph{Derivazione.} Ogni coppia di vertici determina un segmento: $C(10,2) = 45$. Di questi, 10 sono i lati del poligono e non diagonali, quindi $45 - 10 = 35$.

\paragraph{Distrattori.}
\begin{enumerate}
  \item $45$ --- \texttt{sides-included} (\texttt{constraint-ignored}): Hai contato anche i lati: ogni coppia di vertici consecutivi dà un lato, non una diagonale; vanno tolti gli 10 lati.
  \item $90$ --- \texttt{ordered-pairs} (\texttt{order-irrelevant}): Hai contato le coppie ordinate di vertici: il segmento da A a B è lo stesso che va da B ad A, e anche i lati vanno tolti.
  \item $7$ --- \texttt{one-vertex-only} (\texttt{other}): Hai contato solo le diagonali che partono da un vertice ($10 - 3$): ognuno dei 10 vertici ne è capolinea, e ogni diagonale ha due estremi.
\end{enumerate}

\section{Urne}

\subsection{\texorpdfstring{\texttt{urna-senza-reimmissione-entrambe-rosse} (intermedio)}{urna-senza-reimmissione-entrambe-rosse (intermedio)}}

\paragraph{Quesito.} Un'urna contiene 6 palline rosse e 2 palline blu. Si estraggono due palline in sequenza senza reimmissione. Qual è la probabilità che entrambe siano rosse?

\paragraph{Risposta corretta.} $15/28$

\paragraph{Derivazione.} La prima estrazione ha probabilità $6/8$. La seconda, sapendo che è rimasta una rossa in meno, ha probabilità $5/7$. Il prodotto è $6/8 \cdot  5/7 = 15/28$. Usare $6/8$ anche alla seconda estrazione ignora che l'urna è cambiata.

\paragraph{Distrattori.}
\begin{enumerate}
  \item $3/14$ --- \texttt{second-draw-blue} (\texttt{other}): Alla seconda estrazione hai usato la probabilità di una blu ($2/7$): dopo una rossa restano 5 rosse su 7.
  \item $9/16$ --- \texttt{with-replacement} (\texttt{repetition}): Hai trattato le estrazioni come indipendenti (con reimmissione): senza reimmissione, alla seconda estrazione il totale e il numero di rosse diminuiscono.
  \item $3/4$ --- \texttt{first-draw-only} (\texttt{constraint-applied-once}): È la probabilità della sola prima estrazione: serve anche che la seconda sia rossa, e l'urna nel frattempo è cambiata.
\end{enumerate}

\subsection{\texorpdfstring{\texttt{urna-senza-reimmissione-stesso-colore} (intermedio)}{urna-senza-reimmissione-stesso-colore (intermedio)}}

\paragraph{Quesito.} Un'urna contiene 5 palline rosse e 7 palline blu. Si estraggono due palline contemporaneamente. Qual è la probabilità che abbiano lo stesso colore?

\paragraph{Risposta corretta.} $31/66$

\paragraph{Derivazione.} I casi possibili sono $C(12,2) = 66$ coppie. I casi favorevoli sono le coppie tutte rosse più quelle tutte blu: $C(5,2) + C(7,2) = 10 + 21$. Quindi $P = 31/66 = 31/66$.

\paragraph{Distrattori.}
\begin{enumerate}
  \item $37/72$ --- \texttt{with-replacement} (\texttt{repetition}): Hai trattato le due estrazioni come indipendenti (con reimmissione): estraendo contemporaneamente la prima pallina non torna nell'urna, quindi le coppie vanno contate senza ripetizioni.
  \item $31/33$ --- \texttt{ordered-over-unordered} (\texttt{order-irrelevant}): Hai contato i casi favorevoli come coppie ordinate ($5\cdot 4 + 7\cdot 6$) ma i casi possibili come coppie non ordinate $C(12,2)$: numeratore e denominatore devono usare lo stesso modello.
  \item $5/33$ --- \texttt{red-pairs-only} (\texttt{additive-union}): Hai contato solo le coppie tutte rosse: \guillemotleft{}stesso colore\guillemotright{} comprende anche le coppie tutte blu, che vanno sommate.
\end{enumerate}

\subsection{\texorpdfstring{\texttt{urna-almeno-uno} (avanzato)}{urna-almeno-uno (avanzato)}}

\paragraph{Quesito.} Un'urna contiene 3 palline rosse e 6 palline blu. Si estraggono due palline contemporaneamente. Qual è la probabilità che sia presente almeno una pallina rossa?

\paragraph{Risposta corretta.} $7/12$

\paragraph{Derivazione.} È più semplice usare l'evento complementare "nessuna rossa", cioè due blu: $C(6,2)/C(9,2) = 15/36 = 5/12$. Quindi $P(almeno una rossa) = 1 - 5/12 = 7/12$.

\paragraph{Distrattori.}
\begin{enumerate}
  \item $1/2$ --- \texttt{exactly-one} (\texttt{constraint-ignored}): Hai calcolato la probabilità di esattamente una rossa ($3\cdot 6/36$): \guillemotleft{}almeno una\guillemotright{} comprende anche il caso di due rosse.
  \item $5/12$ --- \texttt{complement-itself} (\texttt{complement-ignored}): È la probabilità dell'evento contrario (nessuna rossa, cioè due blu): il risultato richiede 1 $-$ quel valore.
  \item $5/9$ --- \texttt{complement-with-replacement} (\texttt{repetition}): Hai calcolato il complementare come se le estrazioni fossero con reimmissione, 1 $-$ ($6/9$)$^{2}$: estraendo insieme, la seconda pallina viene da un'urna con una pallina in meno.
\end{enumerate}

\section{Probabilità elementare}

\subsection{\texorpdfstring{\texttt{probabilita-dadi-somma} (base)}{probabilita-dadi-somma (base)}}

\paragraph{Quesito.} Si lanciano due dadi equi a sei facce. Qual è la probabilità che la somma dei punteggi sia 4?

\paragraph{Risposta corretta.} $1/12$

\paragraph{Derivazione.} I casi possibili sono $6 \cdot  6 = 36$ coppie ordinate. Le coppie con somma 4 sono (1, 3), (2, 2), (3, 1): 3 casi favorevoli. Quindi $P = 3/36 = 1/12$.

\paragraph{Distrattori.}
\begin{enumerate}
  \item $1/9$ --- \texttt{double-counted-twice} (\texttt{order-irrelevant}): Hai contato due volte la coppia doppia (2, 2): scambiando i due dadi si ottiene la stessa coppia, quindi compare una volta sola fra le 36.
  \item $1/11$ --- \texttt{sums-equiprobable} (\texttt{other}): Hai trattato le 11 somme possibili (da 2 a 12) come equiprobabili: le somme centrali si ottengono con più coppie, quindi vanno contate le 36 coppie ordinate.
  \item $2/21$ --- \texttt{unordered-pairs} (\texttt{order-irrelevant}): Hai usato le coppie non ordinate sia al numeratore (2) sia al denominatore (21): le 21 coppie non ordinate non sono equiprobabili, perché (1, 2) può uscire in due modi e (1, 1) in uno solo.
\end{enumerate}

\subsection{\texorpdfstring{\texttt{probabilita-unione-eventi} (intermedio)}{probabilita-unione-eventi (intermedio)}}

\paragraph{Quesito.} Da un mazzo di 52 carte si estrae una carta. Qual è la probabilità che sia una figura (fante, donna o re) oppure di picche?

\paragraph{Risposta corretta.} $11/26$

\paragraph{Derivazione.} $P(A \cup  B) = P(A) + P(B) - P(A \cap  B)$: $12/52 + 13/52 - 3/52$, perché le figure di picche sono 3. Risultato $22/52 = 11/26$. Dimenticare di togliere l'intersezione porterebbe a $25/52$, contando due volte le carte che sono una figura (fante, donna o re) e di picche insieme.

\paragraph{Distrattori.}
\begin{enumerate}
  \item $25/52$ --- \texttt{intersection-not-removed} (\texttt{additive-union}): Hai sommato le due probabilità senza togliere l'intersezione: le carte che sono una figura (fante, donna o re) e di picche (3 in tutto) verrebbero contate due volte.
  \item $1/4$ --- \texttt{spades-only} (\texttt{constraint-ignored}): Hai considerato solo l'evento \guillemotleft{}picche\guillemotright{}, dimenticando la prima condizione del quesito.
  \item $9/52$ --- \texttt{first-event-minus-intersection} (\texttt{inclusion-exclusion-truncated}): Hai contato solo le figure che non sono di picche (9): mancano tutte le carte di picche.
\end{enumerate}

\subsection{\texorpdfstring{\texttt{probabilita-almeno-uno} (intermedio)}{probabilita-almeno-uno (intermedio)}}

\paragraph{Quesito.} Un dado equo a sei facce viene lanciato 2 volte. Qual è la probabilità di ottenere almeno un sei?

\paragraph{Risposta corretta.} $11/36$

\paragraph{Derivazione.} Ogni tentativo fallisce con probabilità $q = 1 - 1/6 = 5/6$. Tutti i 2 tentativi falliscono con probabilità $q^{2} = 25/36$. Quindi $P(almeno uno) = 1 - 25/36 = 11/36$.

\paragraph{Distrattori.}
\begin{enumerate}
  \item $1/3$ --- \texttt{added-attempts} (\texttt{additive-union}): Hai sommato le probabilità dei singoli tentativi ($2 \times  1/6$): così conti più volte i casi in cui più tentativi riescono insieme. Il metodo corretto passa dal complementare $1 - q^{n}$.
  \item $25/36$ --- \texttt{complement-itself} (\texttt{complement-ignored}): È la probabilità dell'evento contrario, cioè che nessun tentativo riesca: \guillemotleft{}almeno uno\guillemotright{} si ottiene con $1 - q^{n}$.
  \item $5/36$ --- \texttt{success-in-fixed-position} (\texttt{order-irrelevant}): È la probabilità di un successo in una posizione precisa e insuccesso negli altri 1 tentativi, p $\cdot$ ($1 - p$)\textasciicircum{}1: \guillemotleft{}almeno uno\guillemotright{} ammette successi in qualsiasi posizione e anche più di uno.
\end{enumerate}

\section{Probabilità avanzata}

\subsection{\texorpdfstring{\texttt{probabilita-totale-bayes} (avanzato)}{probabilita-totale-bayes (avanzato)}}

\paragraph{Quesito.} In una fabbrica la macchina 1 produce 240 pezzi su 400 e ne difetta il $1/24$, mentre la macchina 2 ne produce 160 su 400 con tasso di difetti $1/8$. Si estrae a caso un pezzo e si scopre che è difettoso. Qual è la probabilità che provenga dalla macchina 1?

\paragraph{Risposta corretta.} $1/3$

\paragraph{Derivazione.} Fra i 30 pezzi difettosi, 10 vengono dalla macchina 1 e 20 dalla macchina 2. Quindi $P(M_{1}\mid difettoso) = 10/30 = 1/3$. È il teorema di Bayes: $P(M_{1}\mid D) = P(D\mid M_{1})\cdot P(M_{1}) / P(D)$.

\paragraph{Distrattori.}
\begin{enumerate}
  \item $3/5$ --- \texttt{prior-share} (\texttt{other}): È la probabilità a priori che un pezzo venga dalla macchina 1: la condizione \guillemotleft{}sapendo che è difettoso\guillemotright{} restringe il campione ai soli difettosi, e in quel gruppo le proporzioni cambiano.
  \item $1/24$ --- \texttt{conditioning-inverted} (\texttt{conditioning-inverted}): Hai confuso $P(macchina 1 \mid  difettoso)$ con il tasso di difetti $P(difettoso \mid  macchina 1)$: il teorema di Bayes inverte proprio la direzione del condizionamento.
  \item $2/3$ --- \texttt{machine2-share} (\texttt{complement-ignored}): Hai calcolato la quota della macchina 2 fra i difettosi: per la macchina 1 basta prenderne il complementare.
\end{enumerate}

\subsection{\texorpdfstring{\texttt{bernoulli-esattamente-k} (avanzato)}{bernoulli-esattamente-k (avanzato)}}

\paragraph{Quesito.} Una moneta equa viene lanciata 3 volte. Qual è la probabilità di ottenere esattamente 1 teste?

\paragraph{Risposta corretta.} $3/8$

\paragraph{Derivazione.} Schema di Bernoulli con $n = 3$, $k = 1$, $p = 1/2$. Il coefficiente binomiale è $C(3,1) = 3$, perché conta le sequenze con esattamente 1 teste. La probabilità di ciascuna sequenza è ($1/2$)\textasciicircum{}3. Quindi $P = C(3,1)$ $\cdot$ ($1/2$)\textasciicircum{}$3 = 3/8 = 3/8$.

\paragraph{Distrattori.}
\begin{enumerate}
  \item $1/8$ --- \texttt{single-sequence} (\texttt{order-irrelevant}): Hai contato una sola sequenza: esistono $C(3,1) = 3$ sequenze con esattamente 1 teste, e vanno tutte contate.
  \item $1/4$ --- \texttt{two-tosses} (\texttt{other}): Hai calcolato la probabilità di un successo e un insuccesso in due lanci: qui i lanci sono 3 e i successi esattamente 1.
  \item $3/16$ --- \texttt{wrong-denominator} (\texttt{other}): Hai diviso per $2^{4}$: i lanci sono 3, quindi le sequenze possibili sono $2^{3}$.
\end{enumerate}

\subsection{\texorpdfstring{\texttt{principio-cassetti-dirichlet} (intermedio)}{principio-cassetti-dirichlet (intermedio)}}

\paragraph{Quesito.} Un mazzo è diviso in 4 semi diversi, con tante carte per ciascun seme. Quante carte occorre estrarre dal mazzo mescolato per essere certi di averne almeno 2 dello stesso seme?

\paragraph{Risposta corretta.} $5$

\paragraph{Derivazione.} Caso peggiore: le prime 4 estrazioni distribuiscono al massimo 1 pezzi per ciascuno dei 4 colori. Il pezzo successivo, il numero 5, non ha scelta: completa il 2-esimo esemplare di qualche colore. Quindi bastano 5 estrazioni.

\paragraph{Distrattori.}
\begin{enumerate}
  \item $2$ --- \texttt{best-case} (\texttt{other}): È il minimo sperabile: 2 estrazioni possono anche risultare tutte di colori diversi. La certezza si ottiene solo ragionando sul caso peggiore del principio dei cassetti.
  \item $4$ --- \texttt{last-draw-missing} (\texttt{other}): Ci manca l'ultima estrazione: con 4 pezzi puoi ancora avere solo 1 esemplari per ciascun colore; quello successivo crea per forza il 2-esimo di un colore.
  \item $8$ --- \texttt{k-of-every-colour} (\texttt{other}): Hai contato 2 pezzi per ogni colore ($4 \cdot  2$): basta che un solo colore arrivi a 2, non tutti.
\end{enumerate}

\section{Archetipi da Trombetta}

\subsection{\texorpdfstring{\texttt{distribuzione-oggetti-stelle-barre} (avanzato)}{distribuzione-oggetti-stelle-barre (avanzato)}}

\paragraph{Quesito.} In quanti modi distinti si possono distribuire 5 caramelle identiche fra 3 bambini, senza alcun vincolo (qualcuno può restare a bocca asciutta)?

\paragraph{Risposta corretta.} $21$

\paragraph{Derivazione.} Oggetti identici e destinatari distinti: è il modello \guillemotleft{}stelle e barre\guillemotright{}. Disponiamo 5 stelle e 2 barre separatrici, e scegliamo dove mettere le barre: $C(5 + 3 - 1, 3 - 1) = C(7, 2) = 21$.

\paragraph{Distrattori.}
\begin{enumerate}
  \item $243$ --- \texttt{objects-distinguished} (\texttt{labelled-groups}): Hai trattato caramelle identiche come distinguibili: $3^{5}$ conta le assegnazioni di oggetti diversi. Qui gli oggetti sono identici, quindi conta solo quanti ne riceve ciascuno.
  \item $10$ --- \texttt{simple-combinations} (\texttt{repetition}): Hai usato le combinazioni semplici $C(5, 2)$: qui le ripetizioni sono ammesse (uno stesso destinatario può ricevere più oggetti), quindi serve la formula delle combinazioni con ripetizione $C(k + n - 1, n - 1)$.
  \item $6$ --- \texttt{minimum-one-imposed} (\texttt{other}): Hai imposto che ciascuno riceva almeno un oggetto, $C(5 - 1, 3 - 1)$: il testo non pone alcun minimo, quindi qualcuno può non ricevere nulla.
\end{enumerate}

\subsection{\texorpdfstring{\texttt{permutazioni-cicliche} (avanzato)}{permutazioni-cicliche (avanzato)}}

\paragraph{Quesito.} In quanti modi distinti si possono far sedere 7 persone attorno a un tavolo rotondo, se due di loro non devono sedersi vicine e le sistemazioni che differiscono per una rotazione si considerano uguali?

\paragraph{Risposta corretta.} $480$

\paragraph{Derivazione.} Fissiamo la prima persona speciale per eliminare le rotazioni. La seconda non può occupare i due posti adiacenti, quindi ha $7 - 3 = 4$ scelte; le altre 5 persone si dispongono in $(7 - 2)!$ modi. Totale: $4 \cdot  120 = 480$.

\paragraph{Distrattori.}
\begin{enumerate}
  \item $720$ --- \texttt{constraint-ignored} (\texttt{constraint-ignored}): Hai contato tutte le sistemazioni circolari senza applicare il vincolo: vanno escluse quelle in cui le due persone sono adiacenti.
  \item $240$ --- \texttt{excluded-cases} (\texttt{complement-ignored}): Hai contato proprio i casi da escludere, cioè quelli con le due persone vicine: il risultato cercato è il complementare.
  \item $600$ --- \texttt{one-neighbour-excluded} (\texttt{constraint-applied-once}): Hai contato 5 posti liberi per la seconda persona: dopo aver fissato la prima, i posti da evitare sono i due adiacenti, quindi ne restano 4.
\end{enumerate}

\subsection{\texorpdfstring{\texttt{binomio-newton} (avanzato)}{binomio-newton (avanzato)}}

\paragraph{Quesito.} Nello sviluppo di ($x + 3$)\textasciicircum{}6 secondo la formula del binomio di Newton, quale è il coefficiente numerico del termine in $x^{2}$?

\paragraph{Risposta corretta.} $1215$

\paragraph{Derivazione.} La formula del binomio dà il termine generale $C(6,k)\cdot$(x)\textasciicircum{}k$\cdot$(3)\textasciicircum{}($6-k$). Per $k = 2$: $C(6,2) = 15$, $1^{2} = 1$ e $3^{4} = 81$. Coefficiente = $15 \cdot  1 \cdot  81 = 1215$.

\paragraph{Distrattori.}
\begin{enumerate}
  \item $15$ --- \texttt{coefficient-only} (\texttt{other}): Hai preso solo il coefficiente binomiale $C(6,2) = 15$: va moltiplicato anche per le potenze delle basi, $1^{2}$ e $3^{4}$.
  \item $135$ --- \texttt{exponents-swapped} (\texttt{other}): Hai scambiato gli esponenti: nel termine in $x^{2}$ la potenza 2 spetta al primo addendo e la potenza 4 al secondo, non il contrario.
  \item $2430$ --- \texttt{dispositions-used} (\texttt{order-irrelevant}): Hai usato le disposizioni $6!/(6-2)!$ al posto del coefficiente binomiale: nel binomio di Newton serve $C(6,2) = 15$, che divide anche per $2!$.
\end{enumerate}

\subsection{\texorpdfstring{\texttt{inclusione-esclusione} (intermedio)}{inclusione-esclusione (intermedio)}}

\paragraph{Quesito.} In una scuola di 150 studenti, 46 praticano il nuoto e 48 praticano l'atletica; 19 praticano entrambe le discipline. Quanti studenti praticano almeno una delle due discipline?

\paragraph{Risposta corretta.} $75$

\paragraph{Derivazione.} Principio di inclusione-esclusione per due insiemi: |A $\cup$ B| = |A| + |B| $-$ |A $\cap$ B| = $46 + 48 - 19 = 75$.

\paragraph{Distrattori.}
\begin{enumerate}
  \item $94$ --- \texttt{intersection-not-removed} (\texttt{additive-union}): Hai sommato i due gruppi senza togliere l'intersezione: i 19 studenti che praticano entrambe le discipline vengono così contati due volte.
  \item $56$ --- \texttt{intersection-removed-twice} (\texttt{inclusion-exclusion-truncated}): Hai sottratto l'intersezione due volte: chi pratica entrambe le discipline va contato una volta, quindi si toglie 19 una sola volta.
  \item $131$ --- \texttt{school-total} (\texttt{other}): Hai lavorato sul totale della scuola: la domanda riguarda solo chi pratica almeno una disciplina, non chi resta escluso.
\end{enumerate}

